
\begin{proposition}
\label{XIII.1.6}
Sean $f\colon X\to Y$ y $g\colon Y\to Z$ dos $S$-morfismos,
y sea~$\Phi$ un stack sobre~$X$.
\begin{enumerate}
\item[1)] Supongamos que $(\Phi,f)$ y $(f_*\Phi,g)$ son
cohomológicamente propios relativamente a~$S$ en dimensión
$\leqslant -1$ (\resp $\leqslant 0$, \resp $\leqslant 1$). Entonces
lo mismo ocurre con~$(\Phi,gf)$.
\item[2)] Supongamos que $(\Phi,gf)$ es cohomológicamente propio
relativamente a~$S$ en dimensión $\leqslant -1$ (\resp que $(\Phi,gf)$
es cohomológicamente propio relativamente a~$S$ en dimensión
$\leqslant 0$ y $(\Phi,f)$ cohomológicamente propio relativamente
a~$S$ en dimensión $\leqslant -1$, \resp que $(\Phi,gf)$ es
cohomológicamente propio relativamente a~$S$ en dimensión
$\leqslant 1$ y $(\Phi,f)$ cohomológicamente propio relativamente
a~$S$ en dimensión $\leqslant 0$). Entonces $(f_*\Phi,g)$ es
cohomológicamente propio relativamente a~$S$ en dimensión
$\leqslant -1$ (\resp en dimensión $\leqslant 0$, \resp en dimensión
$\leqslant 1$).
\end{enumerate}
\end{proposition}

Para todo $S$-esquema~$S'$, considérese el diagrama siguiente, en el
que todos los cuadrados son cartesianos:
\begingroup
\renewcommand{\thesubsection}{\theproposition}% La etiqueta fuente pertenece al ítem XIII.1.6.
\begin{equation*}
\label{eq:XIII.1.6.1}
\tag{\thesubsection.1}
\begin{array}{c}
\xymatrix{ X' \ar[r]^{f'} \ar[d]_h & Y'
\ar[r]^{g'} \ar[d]_k & Z' \ar[r] \ar[d]_{m} & S' \ar[d] \\ X \ar[r]^f
& Y \ar[r]^g & Z \ar[r] & \,S. }
\end{array}
\end{equation*}
\endgroup
\ifthenelse{\boolean{orig}}
{el\space}
{El\space}%
morfismo canónico
$$
m^*(g_*f_*\Phi)\to g'_*f'_*(h^*\Phi)
$$
se identifica
\marginpar{355}
con el compuesto de los morfismos canónicos
$$
m^*(g_*f_*\Phi)\lto{i}g'_*(k^*f_*\Phi)
\lto{j}g'_*f'_*(h^*\Phi).
$$
\begin{enumerate}
\item[1)] La hipótesis implica que~$i$ y~$j$ son fieles
(\resp plenamente fieles, \resp equivalencias); por tanto, lo mismo
ocurre con~$ji$.
\item[2)] La hipótesis implica que~$ji$ es fiel (\resp que~$ji$ es
\ifthenelse{\boolean{orig}}
{plenamnete\space}
{plenamente\space}%
fiel y~$j$ fiel, \resp que~$ji$ es una equivalencia y~$j$ plenamente
fiel); de ello resulta que~$i$ es fiel (\resp plenamente fiel, \resp una
equivalencia).
\end{enumerate}
